\section{K-032: normalization and graded coefficients of the amplified components}

The result below identifies the entire normalization of an actual collision
component from K-029.  It does not replace its contribution to the global
Hilbert series by the Hilbert series of its normalization.  In particular,
the conductor and the moving Koszul class can affect lower Laurent terms.

\subsection*{The fixed representation and its normalization}

Fix $l\ge1$, and use the K-020 tensor after $k=8l$ amplifications.  Put
\[
 n=2+56l,\qquad m=2+49l,\qquad r=1+7l,\qquad
 M=2^{n-1}=2\cdot256^{7l}.
                                                        \tag{K-032a}
\]
Let $h$ be the diagonal fixing element inducing a primitive grading
element $\xi$ of order $M$.  Let $Y_l$ be the reduced irreducible fixed
component through the quotient of this tensor.  The three spectra of
$h$ each consist of one eigenvalue of multiplicity $r$ and $m-1$ simple
eigenvalues.  Denote the corresponding repeated eigenspaces by
$A,B,C$, each of dimension $r$.

The prefix after $7l$ amplifications has a selected edge $e$.  Write
$\rho$ for its critical margin-one mass.  Its reduced denominator is
exactly $M$, by K-020.  There are two other prefix edges through the
selected first coordinate, one through the selected second coordinate,
and two through the selected third coordinate.  Write their respective
masses as
\[
 \alpha_1,\alpha_2;\qquad \beta;\qquad \gamma_1,\gamma_2.
\]
All are positive, and
\[
 \alpha_1+\alpha_2=\beta=\gamma_1+\gamma_2=1-\rho.
                                                        \tag{K-032b}
\]
Every prefix mass has denominator dividing $M$.

Let
\[
 L=C_{GL_n^3}(h),\qquad C_h=C_{SL_n^3}(h),\qquad
 B_l=\mathbb C[V_n^h]^{C_h}.
\]
Thus $L$ is the product of three groups
$GL_r\times(\mathbb C^*)^{m-1}$, and $C_h$ is the kernel of the
three full-factor determinant characters.  In particular, $C_h$ is
connected and reductive.

\begin{theorem}[K-032, normalization]
The restriction of the tensor quotient map induces a finite birational
morphism
\[
                 \operatorname{Spec}B_l\longrightarrow Y_l.
                                                        \tag{K-032c}
\]
Its source is normal; hence it is the normalization of $Y_l$.
As a representation of the block groups $GL(A)\times GL(B)\times GL(C)$,
the fixed vector space is
\[
 V_n^h=(A\otimes B\otimes C)\oplus A^{\oplus2}\oplus B
                  \oplus C^{\oplus2}\oplus\mathbb C^{3m-8}.
                                                        \tag{K-032d}
\]
The last summands and the five vector summands carry the outside
one-dimensional torus characters specified by their collapsed prefix
edges.  The scalar summands in (K-032d) are therefore not asserted to
be trivial representations of the full group $L$.
\end{theorem}

\begin{proof}
The complete spectral collapse in K-021 proves (K-032d): the selected
edge expands to an $r^3$ block, the five other incident edges each
expand to an $r$-dimensional vector space, and the other $3m-8$ edges
remain one-dimensional.  There are no further zero-weight coordinates.

We give a direct finite-restriction argument in order to retain the
precise group and grading.  This is a special case of the general
finite-restriction phenomenon established by Luna (1975); a modern
statement is Bate--Geranios--Martin, Theorem 1.1.  That general theorem
is not needed as an unproved step in the following argument.

Choose the standard Hermitian forms, for which $h$ is unitary.  We
claim that the $C_h$-nullcone on $V_n^h$ is precisely the intersection
of the $SL_n^3$-nullcone with $V_n^h$.  One inclusion is immediate from
$C_h\subset SL_n^3$.  For the other, suppose that $v\in V_n^h$ is
$C_h$-semistable.  Its $C_h$-orbit closure contains a nonzero closed
$C_h$-orbit, with a representative $w$ whose $C_h$ moment map is zero.
Here one uses the usual complex reductive norm-minimization criterion.
The representative remains in $V_n^h$.

The full special-linear moment map at $w$ is fixed by conjugation by
$h$, because $h$ is unitary and fixes $w$.  Under the invariant
Hermitian identification of Lie algebras and their duals, it therefore
lies in the compact Lie algebra of $C_h$.  Its projection to that Lie
algebra is exactly the $C_h$ moment map, which is zero.  Consequently
the full moment map is zero.  Thus $w$ has a nonzero closed
$SL_n^3$-orbit, and $v$ is $SL_n^3$-semistable.  This proves equality
of the nullcones.

Let $A_l$ be the image of the restriction homomorphism
\[
 \mathbb C[V_n]^{SL_n^3}\longrightarrow\mathbb C[V_n^h].
\]
By K-029, the closure of the fixed-space image is exactly $Y_l$, so
$A_l=\mathbb C[Y_l]$.  Equality of nullcones and surjectivity of the
affine $C_h$ quotient imply
\[
              \sqrt{(A_l)_+B_l}=(B_l)_+.
\]
Since $B_l$ is a finitely generated positively graded algebra, it
follows that $B_l/(A_l)_+B_l$ is finite-dimensional.  Lift a homogeneous
vector-space basis of this quotient.  Induction on degree expresses
every element of $B_l$ as an $A_l$-linear combination of those lifts:
after subtracting the appropriate lifts, the remaining expression
lies in $(A_l)_+B_l$ and has factors of strictly smaller degree on
the $B_l$ side.  Thus $B_l$ is a finite $A_l$-module.

It remains to prove birationality.  K-020 provides a nonempty open
part of $Y_l$ whose points have effectively free closed special-linear
orbits.  Every orbit in the special-linear quotient fibre over such
a point is already closed.  Indeed, a nonclosed orbit would have a
strictly lower-dimensional closed orbit in its boundary, while the
closed effectively free orbit already has the maximal possible
effective group dimension.

Suppose $v,w\in V_n^h$ lie over such a point.  Then $w=g v$ for some
$g\in SL_n^3$.  In the effective group $GL_n^3/S$, where
$S=\{(aI,bI,cI):abc=1\}$, the full stabilizer of $v$ has injective
determinant-product character: its character kernel is the trivial
effective special-linear stabilizer.  Both $h$ and $g^{-1}hg$ fix $v$
and have character $\xi$, so their effective classes coincide.
Consequently $g^{-1}hg=s h$ for some $s\in S$.

In each factor, multiplication by the corresponding scalar from $s$
preserves the spectrum of $h$.  That spectrum has a unique eigenvalue
of multiplicity $r>1$, so the scalar must fix that eigenvalue and
must be one.  Hence $g$ centralizes the actual $h$, and $g\in C_h$.
The finite map (K-032c) therefore has one point in each generic
geometric fibre.  In characteristic zero this makes its function-field
degree one.  Finally, an invariant subring of the normal polynomial
ring $\mathbb C[V_n^h]$ is normal.  The finite birational map is the
normalization, as claimed.
\end{proof}

\subsection*{An exact coefficient formula}

Use the tensor grading: polynomial degree $ns$ is denoted by degree
$s$.  For integer $s\ge0$ divisible by $M$, define partitions
\[
 \begin{split}
 P_j(s)&=(\underbrace{s,\ldots,s}_{r-2},s-j,s\rho+j),\\
 Q(s)&=(\underbrace{s,\ldots,s}_{r-1},s\rho).
 \end{split}                                             \tag{K-032e}
\]
The integers $j$ used below satisfy $j\le s(1-\rho)/2$, so these
are partitions in the displayed order.  Their common size is
$s(r-1+\rho)$.

\begin{theorem}[K-032, graded pieces]
The algebra $B_l$ has no graded pieces except in degrees divisible
by $M$.  In those degrees its dimension is the finite exact sum
\[
 \dim (B_l)_s=
 \sum_{j=0}^{\min(s\alpha_1,s\alpha_2)}
 \sum_{v=0}^{\min(s\gamma_1,s\gamma_2)}
          g\bigl(P_j(s),Q(s),P_v(s)\bigr),\qquad M\mid s.
                                                        \tag{K-032f}
\]
All endpoints in this expression are integers.  The convention at
$s=0$ gives $(B_l)_0=\mathbb C$.
\end{theorem}

\begin{proof}
For computing dimensions one may dualize the whole polynomial
representation and use $\operatorname{Sym}(V_n^h)$: dimensions of
invariants are unchanged.  A degree-$d$ vector fixed by $C_h$ belongs
to a character space for $L/C_h\simeq(\mathbb C^*)^3$.  Write that
character as the product of the three full-factor determinants to
integer powers $s_A,s_B,s_C$.  A scalar in any one full factor acts
on degree $d$ by its $d$th power, and on that determinant character
by its $n s_i$th power.  Therefore
\[
                         d=n s_A=n s_B=n s_C.
\]
The common value $s=d/n$ is a nonnegative integer.  In the original
coordinate representation the characters are dualized; the same
grading and invariant dimensions result.

Separate the fixed representation into its summands indexed by the
old prefix edges.  Let $z_{e'}$ be the symmetric degree assigned to
the summand for an edge $e'$.  Invariance under the outside tori
requires margin $s$ at every simple coordinate.  The scalar center
of each repeated $GL_r$ block requires total degree $rs$ at that
selected coordinate.  The prefix incidence columns are linearly
independent.  Its unique margin-one solution is the positive mass
vector $(p_{e'})$, with $p_e=\rho$.  Consequently the only possible
summand degrees are
\[
 z_{e'}=s p_{e'}\quad(e'\ne e),\qquad
 z_e=s(\rho+r-1).                              \tag{K-032g}
\]
Indeed, adding $s(r-1)$ times the selected edge column adds exactly
the three required selected-coordinate margins.  The degrees in
(K-032g) sum to $ns$, since the prefix masses sum to $m$ and
$n=m+r-1$.

The selected degree in (K-032g) can be integral only if $s\rho$ is
integral.  Since $\rho$ has reduced denominator $M$, this requires
$M\mid s$.  Conversely, for such $s$ all summand degrees are
nonnegative integers, because all prefix mass denominators divide
$M$.  Each one-dimensional summand now contributes one monomial.

Put $q=s(r-1+\rho)$.  Schur--Weyl decomposition of the block gives
\[
 \operatorname{Sym}^{q}(A\otimes B\otimes C)
   =\bigoplus_{\lambda,\mu,\nu\vdash q}
       g(\lambda,\mu,\nu)\,
       S_\lambda A\otimes S_\mu B\otimes S_\nu C,
                                                        \tag{K-032h}
\]
with partitions of length at most $r$.  The two first-factor vector
summands contribute
\[
 \operatorname{Sym}^{s\alpha_1}A\otimes
 \operatorname{Sym}^{s\alpha_2}A
 =\bigoplus_{j=0}^{\min(s\alpha_1,s\alpha_2)}
       S_{(s(1-\rho)-j,j)}A,
                                                        \tag{K-032i}
\]
each with multiplicity one, by the Pieri rule.  The analogous formula
holds in the third factor, while the middle factor contributes
$\operatorname{Sym}^{s\beta}B$.

The determinant character $\det^s$ occurs in
$S_\lambda A\otimes S_\eta A$ precisely when $\lambda$ is the
complement of $\eta$ in the $r$-row rectangle $(s^r)$; its
multiplicity is then one.  Complementing the partitions in
(K-032i) gives $P_j(s)$, and complementing the one-row partition
$(s\beta)$ gives $Q(s)$.  Performing this calculation in all three
factors proves (K-032f).
\end{proof}

\subsection*{What the normalization determines about the leading term}

Put
\[
 d=r^3-3r^2+5r-2=343l^3+14l+1,\qquad
 F_l(z)=\sum_{a\ge0}\dim(B_l)_{Ma}\,z^a,
\]
and define the positive rational Hilbert density
\[
                   e_l=\lim_{z\to1}(1-z)^d F_l(z)>0.
                                                        \tag{K-032j}
\]
Existence, positivity and rationality follow from the Hilbert--Serre
theorem for the positively graded algebra $B_l$.  This definition
does not assume that its degree semigroup is generated by one.

The finite normalization quotient $B_l/A_l$ has support of dimension
at most $d-1$.  Hence $\mathbb C[Y_l]$ has the same coefficient $e_l$
of $(1-t^M)^{-d}$ as its normalization.  At the generic smooth
effectively free point of $Y_l$, the rank of the moving localization
class is
\[
 E_l(\xi)=
 \frac{\det(1-h\mid (\mathfrak{sl}_n^3)_{\rm mov})}
      {\det(1-h\mid (V_n)_{\rm mov})}.                 \tag{K-032k}
\]
We use the point-pullback convention of P-F23 and K-014: its action
on cotangents has the same eigenvalues as the point differential,
and the grading eigenvalue in degree $s$ is $\xi^s$.
Rank determines the top Hilbert coefficient of a coherent class on
an irreducible component.  Thus the component's order-$d$ local
Hilbert contribution has coefficient
\[
                              e_l E_l(\xi).           \tag{K-032l}
\]
All eigenvalues in the two moving spaces are nontrivial roots of
two-power order.  Every factor $1-\omega$ is therefore a unit at
every prime over three.  In particular $E_l(\xi)$ is a three-adic
unit.  A cyclic orbit of three such components has order-$d$
coefficient $3e_lE_l(\xi)$, since factor permutation identifies
their tangent representations and gradings.  Its three-adic
valuation is exactly $1+v_3(e_l)$.

No value or three-adic valuation of $e_l$ is evaluated here.  Formula
(K-032f) identifies the precise remaining density calculation using
ordinary Kronecker coefficients with at most $r$ parts.  It does not
control the lower-dimensional conductor or the lower terms of the
moving coherent class.  These terms can affect the simple Laurent
coefficient at $\xi$, even after the top coefficient is known.
Accordingly this normalization result is not an exact-$M$ pole or a
global noncancellation theorem.

\paragraph{Source scope.}
Luna, D. (1975). Adh\'erences d'orbite et invariants.
\emph{Inventiones Mathematicae, 29}, 231--238.
doi:10.1007/BF01389851.  The original article's publisher page
provides a subscription preview; its complete proof has not been
read in this run.  The attribution to Luna is made through the
modern primary source below.
The known general finite-restriction theorem is restated as
Theorem 1.1, printed p.~1, in Bate, M., Geranios, H., and Martin, B.,
\emph{Orbit closures and invariants}, arXiv:1604.00924v2
(5 December 2018); published as Bate, M., Geranios, H., and
Martin, B. (2019). Orbit closures and invariants.
\emph{Mathematische Zeitschrift, 293}(3--4), 1121--1159.
doi:10.1007/s00209-019-02228-6.
That modern statement and its explicit attribution to Luna have
been read directly.  The specialization above supplies its own
nullcone and finite-module proof, while the birationality and the
graded formula use the specific complete spectra and incidence
certificates of K-020--K-021.
